\section{Vocabulary and numerical settings}\label{sec:vocab}
\emph{Forced response.} Below the Hopf the column does not oscillate on its own;
we quantify its response by the lock-in amplitude at the beat $\Delta f$
(\ifsupp main-text\fi\ Eq.~\eqref{eq:lockin}). There is no autonomous rhythm to entrain in the stable-focus
J-curve, timing, and tACS sweeps.

\emph{Frequency dominance and entrainment.} Above the Hopf the column free-runs
at $f_0$. A $1{:}1$ phase-capture test would follow the unwrapped phase difference
from the drive and require absence of secular drift. Our maps instead report
$|f_{\rm out}-\Delta f|<0.3$~Hz, where $f_{\rm out}$ is the largest non-DC FFT
peak. This is a frequency-dominance screen: a forced component can outgrow a
persisting autonomous one without capturing it. When the components coexist,
instantaneous-phase estimates also require their separation
\cite{lachaux1999}. The maps therefore identify a dominance region, not a
measured Arnold tongue.

\emph{Beat-locked rate and gating.} Gating means that the amplitude of a
persisting rhythm is modulated by the beat without frequency capture
\cite{canolty2010,tort2010}. The QIF statistic used here is the lock-in of the
population rate $r_E$ at $\Delta f$, complemented by the rate folded on beat
phase over whole cycles. It measures beat-locked firing. A raw rate fold does
not extract the gamma amplitude, and sidebands can arise from amplitude or phase
modulation. The above-Hopf QIF panels therefore establish coexistence and mixing
of the beat and gamma rhythms; a gamma-amplitude estimate would be needed to
identify gating specifically.

The protocol for Fig.~\ref{fig:entrain} places the column on the limit-cycle side
(input $p$ on the stable alpha-cycle branch below its upper Hopf), drives it
with $s(t)=\varepsilon\cos\Omega t$, and records the dominant frequency. Grid
widths and the clipping check are specified in the code. As the Hopf is approached,
the autonomous cycle shrinks and the frequency-dominance region broadens in these
simulations. Its width and the stable-focus susceptibility describe different
responses; their increase near onset does not make them the same quantity.

\ifjne
\subsection{Lock-in leakage control}\label{supp:lockin}
Out-of-band components can leak into the estimate, particularly the DC level of the JR voltage readout ($\bar v\!\sim\!8$~mV, far
larger than the demodulated signal). The finite-window projection of a constant onto the quadrature references
need not vanish, so we subtract $\bar v$ explicitly; this makes the metric
accurate down to the $10^{-8}$ level needed for the JR kHz
analysis.

\subsection{JR integration and measurement windows}\label{supp:jr-runs}
We integrate with a fixed-step fourth-order Runge--Kutta (RK4) scheme,
$\Delta t=2\times10^{-4}$~s (50 samples per carrier cycle at $100$~Hz; reduced to
$10^{-4}$~s for carrier sweeps up to $300$~Hz). The kHz-carrier runs of
Fig.~\ref{fig:khz_direct} use $\Delta t=2\times10^{-5}$~s for the carrier sweep to
$4$~kHz (25 samples per carrier cycle at $2$~kHz, 12.5 at $4$~kHz) and
$1.5\times10^{-5}$~s for the $2$~kHz output spectrum. Halving each step changes the
reported lock-in amplitudes by less than $1\%$. The $10$~kHz figures quoted in
\S\ref{sec:khz} are analytic evaluations of Eq.~\eqref{eq:khz}, not simulations.

The resonance sweep uses $10$~s of settling and a $20$~s
measurement window. Long windows reduce transients from weakly damped modes and
average beating when autonomous activity persists. Other runs use the windows
specified in their scripts.
The two-tone $2$~kHz example (Fig.~\ref{fig:khz_twotone}) uses $p=400$~s$^{-1}$,
$C=135$, and $0.5$~mV per tone with $\tau=0.2$~ms. From a zero initial state,
it settles for $4$~s and records $4$~s at $\Delta t=1.5\times10^{-5}$~s,
sampling before each RK4 step. Its amplitude spectra use plain-mean subtraction
and a Hann window. Parameter sweeps advance all $N_{\rm grid}$ combinations simultaneously
in a vectorized $(N_{\rm grid},6)$ state array.

\fi
